%%%%%%%%%%%%%%%%%%%%%%%%%%%%%%%%%%%%%%%%%%%%%%%%%%%%%%%%%%%%%%%%
%
%  Template for homework of Machine Learning.
%%%%%%%%%%%%%%%%%%%%%%%%%%%%%%%%%%%%%%%%%%%%%%%%%%%%%%%%%%%%%%%%

\documentclass[11pt,letter,notitlepage]{ctexart}

\usepackage{fontspec}
\setCJKmainfont{KaiTi}[BoldFont=KaiTi,ItalicFont=KaiTi]

% Mise en page
\usepackage[left=2cm, right=2cm, lines=45, top=0.8in, bottom=0.7in]{geometry}
\usepackage{fancyhdr}
\usepackage{fancybox}
\usepackage{graphicx}
\usepackage{pdfpages}
\usepackage{enumitem}
\usepackage{algorithm}
\usepackage{algorithmic}
\usepackage{array}
\usepackage{booktabs}
\renewcommand{\headrulewidth}{1.5pt}
\renewcommand{\footrulewidth}{1.5pt}
\newcommand\Loadedframemethod{TikZ}
\usepackage[framemethod=\Loadedframemethod]{mdframed}
\usepackage{amssymb,amsmath}
\usepackage{amsthm}
\usepackage{thmtools}
\usepackage{bm}
\usepackage{multirow}
\usepackage{makecell}
\usepackage{tabularx}
\usepackage{longtable}
\usepackage{float}
\usepackage{caption}
\usepackage[edges]{forest}

\setlength{\topmargin}{0pt}
\setlength{\textheight}{9in}
\setlength{\headheight}{0pt}
\setlength{\oddsidemargin}{0.25in}
\setlength{\textwidth}{6in}

%%%%%%%%%%%%%%%%%%%%%%%%
%%%%%% Define math operator %%%%%
%%%%%%%%%%%%%%%%%%%%%%%%
\DeclareMathOperator*{\argmin}{\bf argmin}
\DeclareMathOperator*{\argmax}{\bf argmax}
\DeclareMathOperator*{\relint}{\bf relint\,}
\DeclareMathOperator*{\dom}{\bf dom\,}
\DeclareMathOperator*{\intp}{\bf int\,}

\DeclareMathOperator*{\cl}{\bf cl\,}
\DeclareMathOperator*{\bd}{\bf bd\,}
\DeclareMathOperator*{\conv}{\bf conv\,}
\DeclareMathOperator*{\epi}{\bf epi\,}

%%%%%%%%%%%%%%%%%%%%%%%
\pagestyle{fancy}

%%%%%%%%%%%%%%%%%%%%%%%%
%% Define the Exercise environment %%
%%%%%%%%%%%%%%%%%%%%%%%%
\mdtheorem[
topline=false,
rightline=false,
leftline=false,
bottomline=false,
leftmargin=-10,
rightmargin=-10
]{exercise}{\textbf{Exercise}}

%%%%%%%%%%%%%%%%%%%%%%%%
%% Define the Problem environment %%
%%%%%%%%%%%%%%%%%%%%%%%%
\mdtheorem[
topline=false,
rightline=false,
leftline=false,
bottomline=false,
leftmargin=-10,
rightmargin=-10
]{problem}{\textbf{Problem}}

%%%%%%%%%%%%%%%%%%%%%%%
%% Define the Solution Environment %%
%%%%%%%%%%%%%%%%%%%%%%%
\declaretheoremstyle
[
spaceabove=0pt,
spacebelow=0pt,
headfont=\normalfont\bfseries,
notefont=\mdseries,
notebraces={(}{)},
headpunct={:},
headindent={},
postheadspace={ },
bodyfont=\normalfont,
qed=$\blacksquare$,
preheadhook={\begin{mdframed}[style=myframedstyle]},
postfoothook=\end{mdframed},
]{mystyle}

\declaretheorem[style=mystyle,title=Solution]{solution}

\mdfdefinestyle{myframedstyle}{%
    topline=false,
    rightline=false,
    leftline=false,
    bottomline=false,
    skipabove=-6ex,
    leftmargin=-10,
    rightmargin=-10,
    backgroundcolor=black!5}

%%%%%%%%%%%%%%%%%%%%%%%
%% Useful commands
%%%%%%%%%%%%%%%%%%%%%%%
\renewcommand{\eqref}[1]{Eq.~(\ref{#1})}
\newcommand{\figref}[1]{Fig.~\ref{#1}}
\newcommand{\proj}[2]{\textbf{P}_{#2} (#1)}
\newcommand{\lspan}[1]{\textbf{span} (#1)}
\newcommand{\rank}[1]{\textbf{rank} (#1)}
\newcommand{\tr}{\operatorname{tr}}
\newcommand{\RNum}[1]{\uppercase\expandafter{\romannumeral #1\relax}}
\newcommand{\mb}[1]{\mathbf{#1}}

% Definition environment
\theoremstyle{definition}
\newtheorem{definition}{Definition}

%%%%%%%%%%%%%%%%%%%%%%%%%%%%%%%%%%%%%%%%%%
%% Homework info
%%%%%%%%%%%%%%%%%%%%%%%%%%%%%%%%%%%%%%%%%%
\newcommand{\lec}{\text{肖力}}
\newcommand{\posted}{\text{2026.4.14}}
\newcommand{\due}{\text{2026.4.26}}
\newcommand{\hwno}{\text{3}}

%%%%%%%%%%%%%%%%%%%%%%%%%%%%%%%%%%%%
%% Put your information here
%%%%%%%%%%%%%%%%%%%%%%%%%%%%%%%%%%%%
\newcommand{\name}{\text{陈璟皓}}
\newcommand{\id}{\text{PB23010359}}

\lhead{\textbf{\name}}
\rhead{\textbf{\id}}
\chead{\textbf{Homework \hwno}}

\begin{document}
    \vspace*{-4\baselineskip}
    \thispagestyle{empty}

    \begin{center}
        {\bf\large 机器学习B}
    \end{center}

    \noindent
    Lecturer: \lec
    \hfill
    Homework \hwno
    \\
    Posted: \posted
    \hfill
    Due: \due
    \\
    Name: \name
    \hfill
    ID: \id

    \noindent
    \rule{\textwidth}{2pt}

    \medskip

    %%%%%%%%%%%%%%%%%%%%%%%%%%%%%%%%%%%%%%%%%%%%%%%%%%%%%%%%%%%%
    % Exercise 1
    %%%%%%%%%%%%%%%%%%%%%%%%%%%%%%%%%%%%%%%%%%%%%%%%%%%%%%%%%%%%
    \begin{exercise}[基于 LIBSVM 的线性核与高斯核 SVM 比较]
课本习题6.2，即：

试使用 LIBSVM，在西瓜数据集 3.0$\alpha$ 上分别用线性核和高斯核训练一个 SVM，并比较其支持向量的差别。

\begin{table}[H]
\centering
\begin{tabular}{p{0.2\textwidth}p{0.2\textwidth}p{0.2\textwidth}p{0.2\textwidth}}
\hline
编号 & 密度 & 含糖率 & 好瓜 \\
\hline
1  & 0.697 & 0.460 & 是 \\
2  & 0.774 & 0.376 & 是 \\
3  & 0.634 & 0.264 & 是 \\
4  & 0.608 & 0.318 & 是 \\
5  & 0.556 & 0.215 & 是 \\
6  & 0.403 & 0.237 & 是 \\
7  & 0.481 & 0.149 & 是 \\
8  & 0.437 & 0.211 & 是 \\
9  & 0.666 & 0.091 & 否 \\
10 & 0.243 & 0.267 & 否 \\
11 & 0.245 & 0.057 & 否 \\
12 & 0.343 & 0.099 & 否 \\
13 & 0.639 & 0.161 & 否 \\
14 & 0.657 & 0.198 & 否 \\
15 & 0.360 & 0.370 & 否 \\
16 & 0.593 & 0.042 & 否 \\
17 & 0.719 & 0.103 & 否 \\
\hline
\end{tabular}
\caption{西瓜数据集 3.0$\alpha$}
\end{table}
    \end{exercise}

    \newpage

    \begin{solution}
        
    \end{solution}

    \newpage

    %%%%%%%%%%%%%%%%%%%%%%%%%%%%%%%%%%%%%%%%%%%%%%%%%%%%%%%%%%%%
    % Exercise 2
    %%%%%%%%%%%%%%%%%%%%%%%%%%%%%%%%%%%%%%%%%%%%%%%%%%%%%%%%%%%%
    \begin{exercise}[线性可分判定与硬间隔支持向量机]
有下面 8 个训练样本（只有两个属性）
\[
x_1=(2,3),\ x_2=(3,2),\ x_3=(4,4),\ x_4=(4,2)
\]
\[
x_5=(6,4),\ x_6=(6,3),\ x_7=(7,2),\ x_8=(8,3)
\]

以及它们对应的二分类的类别标记为
\[
y_1=y_2=y_3=y_4=-1,\qquad y_5=y_6=y_7=y_8=1.
\]

\begin{enumerate}
    \item 判断这 8 个训练样本是否线性可分？如果线性可分，通过硬间隔最大化，求出基于这 8 个训练数据得到的线性支持向量模型（列出具体步骤，其中 SMO 算法求对偶问题的时候，用 SMO 的计算机编程 code 去算）。

    \item 请问哪些训练样本是上面训练所得模型的支持向量？

    \item 通过学出的模型，决定下面 3 个样本的预测类别
    \[
    x_9=(3,4),\quad x_{10}=(7,4),\quad x_{11}=(5,5).
    \]
\end{enumerate}
    \end{exercise}

    \newpage

    \begin{solution}
        
    \end{solution}

    \newpage

    %%%%%%%%%%%%%%%%%%%%%%%%%%%%%%%%%%%%%%%%%%%%%%%%%%%%%%%%%%%%
    % Exercise 3
    %%%%%%%%%%%%%%%%%%%%%%%%%%%%%%%%%%%%%%%%%%%%%%%%%%%%%%%%%%%%
    \begin{exercise}[朴素贝叶斯分类器与新样本预测]
给出下面 7 只猫在三个属性（颜色、毛发、体型）上的值，以及每只猫对应的标记（性格）是“温顺”或“暴躁”。基于这 7 个训练样本，手算求出朴素贝叶斯分类器来预测一只新猫
\[
x=(\text{黑色},\ \text{粗糙},\ \text{小型})
\]
的性格是哪一类？

\begin{center}
\begin{tabular}{p{0.2\textwidth}p{0.2\textwidth}p{0.2\textwidth}p{0.2\textwidth}}
\toprule
\textbf{颜色} & \textbf{毛发} & \textbf{体型} & \textbf{性格} \\
\midrule
棕色 & 粗糙 & 小型 & 温顺 \\
黑色 & 粗糙 & 大型 & 暴躁 \\
黑色 & 光滑 & 大型 & 暴躁 \\
黑色 & 卷曲 & 小型 & 温顺 \\
白色 & 卷曲 & 小型 & 温顺 \\
白色 & 光滑 & 小型 & 暴躁 \\
红色 & 粗糙 & 大型 & 温顺 \\
\midrule
黑色 & 粗糙 & 小型 & ?? \\
\bottomrule
\end{tabular}
\end{center}
    \end{exercise}

\newpage

\begin{solution}
    
\end{solution}

\end{document}
